\chapter{Introducció}

Aquest treball de recerca té com a punt de partida una realitat propera i quotidiana: l’existència de dues xarxes de bicicletes compartides que conviuen en un mateix territori però que, alhora, no estan connectades entre si.
Es tracta del \textbf{Bicing}, que opera a la ciutat de Barcelona, i de l’\textbf{AMBici}, que cobreix l’Hospitalet i altres municipis de l’àrea metropolitana.
Un exemple clar d’aquesta manca de connexió es troba a la zona de la \textbf{Riera Blanca}, on les dues xarxes es toquen territorialment però no permeten a l’usuari continuar el trajecte amb la mateixa bicicleta.
En aquests casos, la persona ha de deixar la bicicleta d’una xarxa i agafar-ne una altra, amb la incomoditat que això suposa.

La \textbf{justificació} del treball rau, per tant, en la necessitat de buscar una solució pràctica a aquesta situació.
Tot i que actualment els dos serveis funcionen de manera separada, el que es planteja aquí és un intent de \textbf{“connectar-los” digitalment} mitjançant un programa capaç d’integrar les seves dades.
Aquest programa no pot eliminar les barreres físiques ni legals entre les dues xarxes, però sí que pot oferir a l’usuari informació útil: quines estacions hi ha disponibles, quines tenen bicicletes lliures, quines tenen espais per aparcar i, sobretot, on es poden trobar punts de transbord entre Bicing i AMBici.

Així, aquest treball no només té un interès acadèmic, sinó també social i pràctic.
En una ciutat on la mobilitat sostenible és cada vegada més important, poder disposar d’una eina que ajudi a utilitzar dues xarxes de bicicletes com si fossin una sola és un pas endavant en la recerca de solucions més eficients.
A més, aquest projecte també em permet desenvolupar competències personals i tècniques: aprendre a programar de manera estructurada amb Python, treballar amb dades en temps real, aplicar càlculs matemàtics, utilitzar Git per controlar les versions i redactar aquesta memòria amb LaTeX i Kile.

En resum, la introducció vol deixar clar que aquest treball respon a una \textbf{necessitat real}, que és millorar la connexió entre el Bicing i l’AMBici, i que alhora serveix com a espai d’aprenentatge per posar en pràctica eines i coneixements propis del món de la informàtica i de la recerca.


Cita bibliogràfica: \cite{TP} \cite{GNU-ca}



